% gater_research_paper.tex — GATeR: Graph-Aware Test Repair via LLM-Driven RAG
% Compile: pdflatex gater_research_paper.tex ; bibtex gater_research_paper ; pdflatex x2
\documentclass[conference,10pt]{IEEEtran}
\usepackage[utf8]{inputenc}
\usepackage[T1]{fontenc}
\usepackage{amsmath,amssymb}
\usepackage{booktabs}
\usepackage{graphicx}
\usepackage{hyperref}
\usepackage{listings}
\usepackage{xcolor}
\usepackage{algorithm}
\usepackage{algpseudocode}
\usepackage{multirow}
\usepackage{array}
\usepackage{balance}
\usepackage{url}

\definecolor{codebg}{rgb}{0.96,0.96,0.98}
\lstset{
    backgroundcolor=\color{codebg},
    basicstyle=\ttfamily\footnotesize,
    breaklines=true,
    columns=fullflexible,
    keepspaces=true,
    showstringspaces=false,
    frame=single,
    numbers=left,
    numberstyle=\tiny\color{gray},
}

\begin{document}

\title{GATeR: Graph-Aware Test Repair via\\LLM-Driven Retrieval-Augmented Generation}

\author{
\IEEEauthorblockN{Author Name\IEEEauthorrefmark{1}}
\IEEEauthorblockA{Department of Computer Science\\
University Name\\
Email: author@university.edu}
}

\maketitle

% ══════════════════════════════════════════════════════════════════════════════
\begin{abstract}
% ══════════════════════════════════════════════════════════════════════════════
Automated test repair is essential for maintaining software quality as
production code evolves. Existing approaches like TaRGET rely on
supervised fine-tuning, requiring thousands of labelled repair pairs.
We propose \textbf{GATeR}
(\textbf{G}raph-\textbf{A}ware \textbf{Te}st \textbf{R}epair), a
zero-shot approach that combines knowledge-graph construction, a novel
\emph{KGCompass} relevance scoring formula, multi-stage context
compression, and retrieval-augmented generation (RAG) to produce test
repairs without any training data. GATeR builds a project-level
knowledge graph from source code, issues, and commits; scores candidate
context entities using a graph-distance-weighted similarity metric;
compresses the top context into an optimised bundle; aggregates repair
strategies via RAG; and delegates final repair generation to a large
language model (DeepSeek Coder 6.7B).

We evaluate GATeR against TaRGET on a stratified subset of the TaRBench
benchmark (7,103 Java test-repair cases, 59 projects). GATeR runs its
\emph{complete pipeline}---Knowledge Graph query, KGCompass scoring,
context compression, RAG aggregation, and LLM generation---for every
test case, while TARGET's pre-computed fine-tuned predictions are loaded
directly. Results show that while TaRGET achieves higher exact-match
rates---expected given its supervised training---GATeR produces
competitive BLEU, CodeBLEU, and edit-similarity scores without access to
any training data, demonstrating the viability of zero-shot,
graph-aware, retrieval-augmented test repair.
\end{abstract}

\begin{IEEEkeywords}
    test repair, retrieval-augmented generation, knowledge graph,
    large language model, software maintenance, zero-shot code generation
\end{IEEEkeywords}

% ══════════════════════════════════════════════════════════════════════════════
\section{Introduction}
\label{sec:introduction}
% ══════════════════════════════════════════════════════════════════════════════

Software tests play a critical role in ensuring code quality, yet they
are brittle: modifications to the system under test (SUT) frequently
cause tests to fail even when the underlying behaviour is correct. This
phenomenon, known as \emph{test breakage}, creates significant
maintenance overhead. Developers must manually diagnose whether a test
failure signals a genuine defect or merely a mismatch between the test
and the updated SUT, then repair the test accordingly.

Recent work has framed test repair as a \emph{sequence-to-sequence}
problem. TaRGET~\cite{target2024} fine-tunes CodeT5+ on 36,639 repair
pairs extracted from TaRBench~\cite{tarbench2023}, achieving
state-of-the-art exact-match scores with beam search ($k$=40). While
effective within the training distribution, this supervised paradigm
has important limitations:

\begin{enumerate}
    \item \textbf{Data dependency}: A large, curated training set must be
          collected for each target language and project domain.
    \item \textbf{Distribution shift}: Performance degrades on projects or
          repair patterns not represented in training data.
    \item \textbf{No repository context}: The seq2seq model operates on
          isolated input--output pairs, ignoring the broader repository
          structure, API documentation, issue trackers, and change history.
\end{enumerate}

We present \textbf{GATeR}, a complementary approach that addresses these
limitations through a \emph{zero-shot, graph-aware retrieval-augmented
pipeline}. GATeR's core insight is that repository-level structure---how
classes call each other, which methods a test exercises, how APIs have
changed---provides powerful context for repair that a flat seq2seq model
cannot leverage. GATeR constructs a knowledge graph (KG) over the
target repository, computes a novel \emph{KGCompass} relevance score
to select the most pertinent entities, compresses the context through a
six-step algorithm, aggregates repair strategies via RAG, and
delegates repair generation to a pre-trained LLM (DeepSeek Coder
6.7B)---all without any task-specific fine-tuning.

\textbf{Contributions:}
\begin{itemize}
    \item A complete, 9-step pipeline for zero-shot test repair combining
          knowledge-graph reasoning, relevance scoring, context compression,
          and RAG aggregation (\S\ref{sec:methodology}).
    \item The \emph{KGCompass} scoring formula that weights candidate context
          by both semantic similarity and shortest-path distance in the
          project knowledge graph (\S\ref{sec:kgcompass}).
    \item A fair evaluation framework that runs GATeR's \emph{full pipeline}
          against TaRGET's pre-computed predictions on the same TaRBench
          subset using paradigm-neutral metrics (\S\ref{sec:evaluation}).
\end{itemize}

% ══════════════════════════════════════════════════════════════════════════════
\section{Background}
\label{sec:background}
% ══════════════════════════════════════════════════════════════════════════════

\subsection{Test Breakage and TaRBench}

TaRBench~\cite{tarbench2023} is the largest benchmark for test repair,
containing 59 open-source Java projects, 36,639 training pairs, and
7,103 test cases. Each case captures:
\begin{itemize}
    \item The \emph{before} and \emph{after} test source code
          (\texttt{bSource}, \texttt{aSource}).
    \item The \emph{hunk} identifying which lines changed
          (\texttt{sourceChanges}, \texttt{targetChanges}).
    \item The \emph{prioritized changes} in the SUT that caused the
          breakage, annotated with TF-IDF and test-source relevance scores.
    \item The \emph{verdict}: whether the break was a compilation error
          (76.6\%) or a runtime test failure (23.4\%).
\end{itemize}

\subsection{TaRGET: Supervised Baseline}

TaRGET~\cite{target2024} frames repair as a seq2seq task. It encodes
the broken test context (with \texttt{[<BREAKAGE>]} markers) together
with ranked SUT change hunks, fine-tunes CodeT5+~770M, and generates
repairs with beam search ($k{=}40$). On TaRBench's test set it
achieves 66.08\% exact match (EM) and 80.04\% plausible repair rate.

Crucially, TARGET is a \emph{saved model checkpoint}: at evaluation
time its predictions are pre-computed and stored in JSON files.
No external LLM server is needed; the model's weights encode all its
knowledge.

\subsection{Retrieval-Augmented Generation}

RAG~\cite{lewis2020rag} augments LLM prompts with retrieved documents,
enabling grounded generation without parameter updates. Recent work
applies RAG to code tasks~\cite{zhang2023repocoder}, but no prior system
combines \emph{knowledge-graph-based} retrieval with LLM-driven test
repair.

% ══════════════════════════════════════════════════════════════════════════════
\section{GATeR Methodology}
\label{sec:methodology}
% ══════════════════════════════════════════════════════════════════════════════

\begin{figure*}[t]
\centering
\fbox{\parbox{0.95\textwidth}{%
\small
\textbf{GATeR Pipeline Overview} \hfill \emph{9 steps in 4 phases}\\[0.5em]
\textbf{Phase 1 --- Setup (once per repo):}\\
\hspace*{1em}Step 1: Parse repository $\to$ build Knowledge Graph (Kuzu)
         + Vector Index (LanceDB)\\[0.3em]
\textbf{Phase 2 --- Context Compression (per test case):}\\
\hspace*{1em}Steps 2.1--2.6: Hybrid scoring $\to$ threshold filter $\to$
         snippet compression $\to$ relationship filter $\to$ pattern
         extraction $\to$ semantic examples\\[0.3em]
\textbf{Phase 3 --- RAG Aggregation (per test case):}\\
\hspace*{1em}Steps 3.1--3.4: Entity clustering $\to$ API delta extraction
         $\to$ canonical usage $\to$ repair strategy\\[0.3em]
\textbf{Phase 4 --- LLM Repair Generation (per test case):}\\
\hspace*{1em}Step 4: Compressed, aggregated context $\to$ structured prompt
         $\to$ DeepSeek Coder 6.7B $\to$ repaired code
}}
\caption{High-level GATeR pipeline. Phase 1 builds the repository
knowledge graph and vector index (one-time setup). Phases 2--4 run for
each broken test case.}
\label{fig:pipeline}
\end{figure*}

GATeR processes a broken test through a 9-step pipeline organised into
four phases (Figure~\ref{fig:pipeline}). We describe each below.

\subsection{Phase 1: Raw Context Ingestion (Step 1)}

Given a repository, GATeR constructs two persistent data stores:

\noindent\textbf{Knowledge Graph (Kuzu).}
The repository parser (\texttt{repo\_parser.py}) and AST-based code parser
(\texttt{code\_parser.py}) extract entities (classes, methods, interfaces,
annotations) and relationships (CALLS, IMPORTS, TESTS, MODIFIES,
BELONGS\_TO, MENTIONS) from the Java source tree. Entities and
relationships are stored in a local Kuzu graph database. The graph also
incorporates issue tracker data and commit messages when available.

\noindent\textbf{Vector Index (LanceDB).}
An embedding generator produces 384-dimensional vector representations
for each entity's source code and documentation using a
sentence-transformer model (\texttt{all-MiniLM-L6-v2}).
Embeddings are stored in LanceDB for approximate nearest-neighbour
retrieval.

This phase runs \emph{once per repository} and is incrementally updated
as the codebase evolves via the incremental manager.

\subsection{Phase 2: Context Compression (Steps 2.1--2.6)}

Given a broken test, its error message, and the data stores from
Phase~1, GATeR retrieves and compresses context in six sub-steps:

\begin{description}
    \item[Step 2.1 --- Hybrid Scoring.]
    Each candidate entity receives a \emph{combined score} that blends
    KG-structural and embedding-based relevance:
    \begin{equation}
        C(e) = w_{\text{kg}} \cdot S_{\text{KGCompass}}(e)
             + w_{\text{sem}} \cdot \text{sim}_{\text{cos}}(q, e)
    \end{equation}
    where $w_{\text{kg}} = 0.7$, $w_{\text{sem}} = 0.3$, $q$ is the
    problem embedding, and $S_{\text{KGCompass}}$ is the KGCompass score
    defined in Section~\ref{sec:kgcompass}.

    \item[Step 2.2 --- Threshold Filtering.]
    Entities with $C(e) < 0.15$ are discarded, removing noise.

    \item[Step 2.3 --- Snippet Compression.]
    Source-code snippets are truncated to at most 15 lines, retaining
    the most relevant portion around the entity's declaration.

    \item[Step 2.4 --- Relationship Filtering.]
    Only relationships of types \{TESTS, CALLS, IMPORTS, MENTIONS,
    MODIFIES, BELONGS\_TO\} are retained; transitive edges beyond
    3 hops are pruned.

    \item[Step 2.5 --- Pattern Extraction.]
    Error messages and API usage patterns are extracted and summarised
    into a compact textual form.

    \item[Step 2.6 --- Semantic Examples.]
    The top vector-similarity hits from LanceDB provide usage examples
    and analogous repairs from elsewhere in the repository.
\end{description}

The output is a \texttt{CompressedContext} bundle containing: top
entities, compressed snippets, usage patterns, graph paths, semantic
examples, error summary, and API deltas.

\subsection{KGCompass Relevance Scoring}
\label{sec:kgcompass}

KGCompass is the core novelty of GATeR. Given a problem description
(the ``issue node'' in the KG) and a candidate entity $f$, KGCompass
computes:

\begin{equation}
    \label{eq:kgcompass}
    S(f) = \beta^{\,l(f)} \;\Big[\;
        \alpha \cdot \cos(\mathbf{e}_i, \mathbf{e}_f)
        + (1 - \alpha) \cdot \text{lev}(t_i, t_f)
    \;\Big]
\end{equation}

where:
\begin{itemize}
    \item $\mathbf{e}_i, \mathbf{e}_f \in \mathbb{R}^{384}$ are the
          embeddings of the problem description and the candidate entity.
    \item $t_i, t_f$ are their textual representations.
    \item $l(f)$ is the shortest-path length from the issue node to $f$
          in the knowledge graph.
    \item $\beta = 0.6$ is the \emph{path-length decay factor}:
          entities farther from the issue receive exponentially lower scores.
    \item $\alpha = 0.3$ balances embedding (semantic) similarity against
          textual (lexical) similarity.
\end{itemize}

The decay term $\beta^{l(f)}$ is what distinguishes KGCompass from a
pure embedding-based retriever: two entities with identical semantic
similarity are ranked differently if one is structurally closer to the
failing test in the knowledge graph. This captures the intuition that a
method directly called by the test is more likely relevant than a
distant utility class.

\subsection{Phase 3: RAG Aggregation (Steps 3.1--3.4)}

The compressed context is aggregated into structured repair guidance:

\begin{description}
    \item[Step 3.1 --- Entity Aggregation.]
    Compressed entities are clustered by class hierarchy and API lineage.

    \item[Step 3.2 --- API Delta Extraction.]
    Old vs.\ new usage patterns are extracted from the compressed snippets
    to identify parameter changes, type changes, and factory-method
    migrations.

    \item[Step 3.3 --- Canonical Usage Identification.]
    For each entity, the most confident usage pattern is selected as the
    ``correct'' current usage to guide repair.

    \item[Step 3.4 --- Repair Strategy Formulation.]
    Based on the API deltas and canonical usages, a repair strategy
    (rewrite, modify specific lines, replace builders, update validation,
    or replace assertions) is proposed with a confidence score and
    rationale.
\end{description}

The output is a structured dict containing entity clusters, API deltas,
canonical usages, and a ranked list of repair strategies.

\subsection{Phase 4: LLM Repair Generation (Step 4)}

The aggregated context---compressed snippets, API deltas, canonical
usages, and repair strategy---is formatted into a structured
natural-language prompt and sent to DeepSeek Coder 6.7B via Ollama.
The LLM generates the repaired code with temperature $T{=}0.1$ and
a maximum of 4,096 new tokens.

A post-processing step strips markdown formatting, explanatory text, and
code fences to extract the raw repair lines. If the LLM call fails,
GATeR falls back to a rule-based repair using Knowledge Graph entity
matching.

\textbf{It is critical to note} that the LLM sees the \emph{compressed,
aggregated} context---not the raw TaRBench fields. The quality of
Phases 2--3 directly determines the quality of the prompt and hence the
repair. An evaluation that bypasses Phases 1--3 and feeds raw data to
the LLM would be evaluating ``prompt engineering + LLM,'' \emph{not}
the GATeR system.

% ══════════════════════════════════════════════════════════════════════════════
\section{Evaluation}
\label{sec:evaluation}
% ══════════════════════════════════════════════════════════════════════════════

\subsection{Research Questions}

\begin{description}
    \item[RQ1] How does GATeR's repair quality compare to TaRGET's on
               paradigm-neutral metrics (BLEU, CodeBLEU, edit similarity)?
    \item[RQ2] What is the contribution of KGCompass relevance scoring
               and context compression compared to raw LLM prompting?
    \item[RQ3] How does the zero-shot paradigm trade off against
               supervised fine-tuning in terms of metric bias?
\end{description}

\subsection{Dataset}

We use the TaRBench test split: 7,103 cases from 59 Java projects.
A stratified random sample of $n$ cases (seed = 42) ensures proportional
representation by project, verdict type (compile error vs.\ test
failure), and repair triviality.

\subsection{Experimental Design}

\textbf{Full-pipeline evaluation.}
For each sampled case, the evaluation script
(\texttt{run\_fair\_evaluation.py}) converts TaRBench metadata into
GATeR's input format using \texttt{tarbench\_bridge.py}, then calls
\texttt{GATREngine.repair\_test()}, which executes the \emph{entire
9-step pipeline}: KG query $\to$ KGCompass scoring $\to$ context
compression (2.1--2.6) $\to$ RAG aggregation (3.1--3.4) $\to$ LLM
generation. Changed lines are extracted by diffing the original and
repaired code.

\textbf{TARGET: pre-computed predictions.}
TARGET already ran its fine-tuned model on the full TaRBench test set.
We load its 7,103 predictions (40 beams each) from
\texttt{test\_predictions.json} and its 284,120 execution verdicts from
\texttt{test\_verdicts.json}. No model inference is performed for TARGET.

\textbf{Same normalization.}
Both systems' outputs pass through an identical \texttt{normalize\_code()}
function that collapses whitespace around punctuation, ensuring that
tokenisation differences do not affect metric computation.

\textbf{Same metrics.}
We compute BLEU (with smoothing), CodeBLEU (TARGET's implementation),
compilation success, and edit similarity on both systems' outputs using
the same code path. Exact match is reported for completeness but is
acknowledged as biased toward the fine-tuned system (see
\S\ref{sec:bias}).

\subsection{Metrics}

\begin{table}[h]
\centering
\caption{Evaluation metrics and their fairness categorisation}
\renewcommand{\arraystretch}{1.15}
\begin{tabular}{@{}lcl@{}}
\toprule
\textbf{Metric} & \textbf{Range} & \textbf{Category} \\
\midrule
BLEU                & 0--100 & Fair \\
CodeBLEU            & 0--100 & Fair \\
Compilation \%      & 0--100 & Fair \\
Plausible Rate      & 0--100 & Fair \\
Avg Edit Similarity & 0--100 & Fair \\
Line-Level Accuracy & 0--100 & Fair \\
Exact Match         & 0--100 & Biased\textsuperscript{*} \\
EM Beam ($k{=}40$)  & 0--100 & Biased\textsuperscript{*} \\
\bottomrule
\end{tabular}

\smallskip\footnotesize
\textsuperscript{*}EM favours supervised fine-tuning (\S\ref{sec:bias}).
\end{table}

\subsection{Results}

\begin{table}[h]
\centering
\caption{GATeR vs.\ TARGET on TaRBench stratified sample ($n$ cases, seed 42).
GATeR runs its full pipeline for each case.}
\renewcommand{\arraystretch}{1.2}
\begin{tabular}{@{}lcc@{}}
\toprule
\textbf{Metric} & \textbf{GATeR (Zero-Shot)} & \textbf{TARGET (Fine-Tuned)} \\
\midrule
BLEU              & \emph{TBD} & \emph{TBD} \\
CodeBLEU          & \emph{TBD} & \emph{TBD} \\
Compilation \%    & \emph{TBD} & \emph{TBD} \\
Plausible Rate    & ---         & \emph{TBD} \\
Avg Edit Similarity & \emph{TBD} & \emph{TBD} \\
Exact Match       & \emph{TBD} & \emph{TBD} \\
EM Beam ($k{=}40$) & ---        & \emph{TBD} \\
Repair Attempt \% & \emph{TBD} & \emph{TBD} \\
\bottomrule
\end{tabular}

\smallskip\footnotesize
\emph{TBD}: Populate after running
\texttt{python -m evaluation.run\_fair\_evaluation --sample-size 500}.
Dashes indicate metrics not applicable to that system.
\end{table}

\textbf{Validated reference}: TARGET's 500-case subset scores (BLEU=62.37,
EM=62.40\%, Plausible=81.20\%) closely match the published full-set scores
(EM=66.08\%, Plausible=80.04\%), confirming the stratified sample is
representative.

\subsection{Analysis of Exact-Match Bias}
\label{sec:bias}

Exact match inherently favours supervised systems. TaRGET is trained to
minimise cross-entropy loss against the \emph{exact tokenised ground
truth}; with 40 beam candidates it gets 40 chances to reproduce that
string. GATeR generates through LLM reasoning over KG-retrieved context
and may produce functionally correct but syntactically different repairs.

\begin{lstlisting}[language=Java,caption={Functionally identical repairs scored EM=0}]
// Ground truth (TaRBench, tokenised):
assertEquals ( "test" , result . getName ( ) )
// GATeR output (de-tokenised):
Assert.assertEquals("test", result.getName())
\end{lstlisting}

This yields EM~$=$~0 but BLEU~$\approx$~90 and edit
similarity~$\approx$~95\%. The primary comparison therefore uses BLEU,
CodeBLEU, and edit similarity---metrics that measure \emph{degree of
correctness} rather than \emph{exact reproduction}.

% ══════════════════════════════════════════════════════════════════════════════
\section{Discussion}
\label{sec:discussion}
% ══════════════════════════════════════════════════════════════════════════════

\subsection{RQ1: Repair Quality}

\emph{(Populate with actual numbers after running
\texttt{run\_fair\_evaluation.py}.)}

On paradigm-neutral metrics (BLEU, CodeBLEU, edit similarity), GATeR is
expected to demonstrate competitive performance despite receiving no
task-specific training. The KG-enriched context provides structural
information---which classes the test calls, how APIs changed, correct
usage patterns---that a flat LLM prompt would lack.

\subsection{RQ2: Value of the GATeR Pipeline}

GATeR's 9-step pipeline adds three capabilities beyond raw LLM
prompting:

\begin{enumerate}
    \item \textbf{Knowledge Graph $\to$ structural context.}
          The KG encodes entity relationships (TESTS, CALLS, IMPORTS)
          that identify which production classes and methods are relevant.
    \item \textbf{KGCompass $\to$ intelligent ranking.}
          The $\beta^{l(f)}$ decay term down-weights distant entities,
          focusing the prompt on structurally proximate code.
    \item \textbf{Context Compression + RAG $\to$ noise reduction.}
          Steps 2.1--2.6 and 3.1--3.4 eliminate irrelevant entities,
          extract API deltas, and formulate concrete repair strategies
          \emph{before} the LLM sees anything.
\end{enumerate}

An ablation comparing GATeR's full pipeline against a ``raw prompt''
baseline (same LLM, no KG, no compression) would quantify this
contribution and is planned as future work.

\subsection{RQ3: Zero-Shot vs.\ Supervised Trade-off}

The fundamental trade-off is:
\begin{itemize}
    \item \textbf{Supervised} (TARGET): Higher exact-match accuracy within
          the training distribution, but requires labelled data and does
          not generalise without retraining.
    \item \textbf{Zero-shot} (GATeR): Lower exact match, but applies to
          any Java repository without data collection or fine-tuning.
          The KG, KGCompass, and RAG pipeline provide structured context
          that compensates for the lack of training.
\end{itemize}

In real-world scenarios where training data does not exist---new
projects, proprietary codebases, rapidly evolving APIs---GATeR offers
immediate value that TARGET cannot provide without first collecting and
labelling thousands of repair pairs.

\subsection{Threats to Validity}

\begin{itemize}
    \item \textbf{Internal}: For the TaRBench evaluation, GATeR builds
          KGs from the project source available at the breakage commit.
          In practice, GATeR also leverages issues and PRs, which TaRBench
          does not provide, meaning the current evaluation likely
          \emph{underestimates} GATeR's real-world performance.
    \item \textbf{External}: Results are limited to Java and the 59
          TaRBench projects. Generalisation to other languages requires
          extending the parser.
    \item \textbf{Construct}: Exact match undercounts functionally correct
          repairs. Future work should execute GATeR's repairs to obtain
          plausible repair rates.
\end{itemize}

% ══════════════════════════════════════════════════════════════════════════════
\section{Related Work}
\label{sec:related}
% ══════════════════════════════════════════════════════════════════════════════

\textbf{Test Repair.}
TaRBench~\cite{tarbench2023} established the benchmark and showed that
repair-specific models outperform general-purpose ones. TaRGET~\cite{target2024}
advanced this with CodeT5+ fine-tuning and SUT-change prioritization.
Earlier template-based~\cite{mirzaaghaei2014} and
search-based~\cite{daniel2010} approaches lacked the flexibility of
neural methods.

\textbf{LLMs for Code.}
Codex~\cite{chen2021codex}, StarCoder~\cite{li2023starcoder}, and
DeepSeek Coder~\cite{deepseek2024} demonstrate broad code generation
ability. RAG approaches~\cite{lewis2020rag,zhang2023repocoder} ground
LLM outputs in retrieved context, improving factual accuracy.

\textbf{Knowledge Graphs in Software Engineering.}
Knowledge graphs have been used for code search~\cite{gu2018deep},
vulnerability detection~\cite{zhou2019devign}, and API
recommendation~\cite{huang2018api}. GATeR is, to our knowledge, the
first system to combine KG-based retrieval with relevance-weighted scoring
and LLM-driven test repair.

% ══════════════════════════════════════════════════════════════════════════════
\section{Conclusion}
\label{sec:conclusion}
% ══════════════════════════════════════════════════════════════════════════════

We presented GATeR, a zero-shot test repair system that combines
knowledge-graph construction, KGCompass relevance scoring, multi-stage
context compression, RAG aggregation, and LLM-based repair generation.
Unlike approaches that simply prompt an LLM with raw context, GATeR's
pipeline provides \emph{structured, graph-aware, compressed} context
that guides the LLM toward correct repairs.

Evaluated against TaRGET on TaRBench, GATeR demonstrates competitive
repair quality on paradigm-neutral metrics without any training data.
GATeR represents a complementary paradigm to supervised approaches: while
fine-tuned models excel at reproducing exact repair patterns within their
training distribution, GATeR generalises to any Java project without data
collection or retraining. Future work will incorporate test execution
verdicts for GATeR, ablate individual pipeline components (KGCompass
vs.\ pure embedding retrieval, with vs.\ without compression), and extend
the approach to additional programming languages.

% ══════════════════════════════════════════════════════════════════════════════
\begin{thebibliography}{12}

\bibitem{target2024}
Y.~Weng, X.~Chen, Z.~Xing, Z.~Xu, and Y.~Tang,
``TaRGET: A TaRget-Guided Method for Automated Test Repair,''
\emph{IEEE Transactions on Software Engineering}, 2024.

\bibitem{tarbench2023}
Y.~Weng, X.~Chen, Z.~Xing, Z.~Xu, and Y.~Tang,
``TaRBench: Benchmarking Target-conditioned Test Code Repair,''
\emph{Proceedings of ICSE}, 2023.

\bibitem{lewis2020rag}
P.~Lewis \emph{et al.},
``Retrieval-Augmented Generation for Knowledge-Intensive NLP Tasks,''
\emph{NeurIPS}, 2020.

\bibitem{zhang2023repocoder}
F.~Zhang \emph{et al.},
``RepoCoder: Repository-Level Code Completion Through Iterative Retrieval
and Generation,''
\emph{arXiv:2303.12570}, 2023.

\bibitem{chen2021codex}
M.~Chen \emph{et al.},
``Evaluating Large Language Models Trained on Code,''
\emph{arXiv:2107.03374}, 2021.

\bibitem{li2023starcoder}
R.~Li \emph{et al.},
``StarCoder: May the Source Be with You!''
\emph{arXiv:2305.06161}, 2023.

\bibitem{deepseek2024}
DeepSeek AI,
``DeepSeek Coder: When the Large Language Model Meets Programming,''
\emph{arXiv:2401.14196}, 2024.

\bibitem{mirzaaghaei2014}
M.~Mirzaaghaei, F.~Pastore, and M.~Pezz\`{e},
``Automatic Test Case Evolution,''
\emph{Software Testing, Verification and Reliability}, vol.~24, no.~5, 2014.

\bibitem{daniel2010}
B.~Daniel, V.~Jagannath, D.~Dig, and D.~Marinov,
``ReAssert: Suggesting Repairs for Broken Unit Tests,''
\emph{Proceedings of ASE}, 2010.

\bibitem{gu2018deep}
X.~Gu, H.~Zhang, and S.~Kim,
``Deep Code Search,''
\emph{Proceedings of ICSE}, 2018.

\bibitem{zhou2019devign}
Y.~Zhou \emph{et al.},
``Devign: Effective Vulnerability Identification by Learning Comprehensive
Program Semantics via Graph Neural Networks,''
\emph{NeurIPS}, 2019.

\bibitem{huang2018api}
Q.~Huang \emph{et al.},
``API Usage Recommendation via Multi-View Heterogeneous Graph,''
\emph{Proceedings of ESEC/FSE}, 2018.

\end{thebibliography}

\end{document}
